% ECONOMICS_PROBLEM: {"id": "F12-265", "title": "Contracting frictions in global sourcing", "jel_code": "F12", "source_id": "OP-0349"}
% !TeX program = xelatex
\documentclass[11pt,a4paper]{article}
\usepackage[margin=23mm]{geometry}
\usepackage{fontspec}
\setmainfont{Latin Modern Roman}
\usepackage{amsmath,amssymb,mathtools}
\usepackage{xcolor,enumitem,booktabs,longtable,array,xurl,hyperref}
\definecolor{muted}{HTML}{53636C}
\hypersetup{colorlinks=true,urlcolor=blue,linkcolor=blue}
\setlength{\parindent}{0pt}
\setlength{\parskip}{5pt}
\newcommand{\R}{\mathbb R}
\newcommand{\E}{\mathbb E}
\newcommand{\N}{\mathbb N}
\newcommand{\ind}{\mathbf 1}
\newcommand{\Var}{\operatorname{Var}}
\newcommand{\Cov}{\operatorname{Cov}}
\newcommand{\supp}{\operatorname{supp}}
\DeclareMathOperator*{\argmin}{arg\,min}
\DeclareMathOperator*{\argmax}{arg\,max}
\newcommand{\entrymeta}[3]{{\sffamily\small\color{muted}\textbf{\texttt{#1}}\quad|\quad #2\par #3\par}}
\newcommand{\Problemref}[1]{\texttt{#1}}
\newcommand{\JELref}[2]{\texttt{#2} (JEL \texttt{#1})}
\begin{document}
\section*{Contracting frictions in global sourcing}
\textbf{Problem ID:} \texttt{F12-265}\quad \textbf{JEL:} \texttt{F12}\par
% BEGIN ECONOMICS STATEMENT
\label{problem:OP-0349}
{\sffamily\small\color{muted}Original subject group: Trade, globalization and geoeconomics\par}
\label{definitions:problem:30007631}
\textbf{Audit classification:} Background; proposed extension. The NBER primary indexed abstract and metadata verify this real 2025 paper. Full text was inaccessible during this audit, and no exact open passage was verified. Its conditional quantitative question should not be relabeled an unresolved general theorem.
\subsection*{Definitions and precise research target}
Fix heterogeneous firms choosing domestic or foreign suppliers and whether to own production stages. Contracts are incomplete in a specified sense: a declared subset of investment or quality actions cannot be enforced by courts. Let \(\kappa_{ab}\) measure enforcement costs or noncontractible exposure for transactions between jurisdictions \(a,b\). Firms choose organization, sourcing and relationship-specific investment optimally, with wages and prices clearing. Let \(W(\kappa)\) be discounted household welfare and \(Q(\kappa)\) trade volumes. The task is to identify losses from these frictions relative to a defined counterfactual \(\kappa=0\):
\[\Delta W=W(0)-W(\kappa),\qquad\Delta Q=Q(0)-Q(\kappa).\]
Specify which institutional change implements friction removal and which technologies, preferences and trade policies remain fixed. Separate enforcement frictions from search, transport and financing costs using a stated exclusion or structural restriction. Establish whether organizational choices and sourcing patterns uniquely identify \(\kappa\) or only composite wedges. Include endogenous multinational ownership, entry and supplier investment in welfare comparisons. Report sensitivity to bargaining rules and cross-country institutional measurement. A model matching observed sourcing shares can estimate a conditional counterfactual without proving that its enforcement mechanism is uniquely identified or that the broader research question was first stated in that paper.

\textbf{Additional formal definitions and scope (editorial).} Specify a complete contracting model: the nonverifiable action set, feasible enforceable contract terms, outside options, allocation of ownership/control rights and any bargaining rule. Noncontractibility is not simply a proportional transaction cost; if \(\kappa_{ab}\) is used, define whether it changes enforcement probability, investment appropriability, or a resource wedge and identify each component separately. The intervention \(\kappa=0\) must say whether it makes all actions verifiable or only removes enforcement costs; ownership responses and transfers may differ. Define trade volume in physical quantities or a common-price real index, and welfare by discounted household utility with fiscal and contractual resource costs. Even if enforcement improves, \(\Delta Q\) need not be positive for every trade flow; no sign theorem is assumed. A structural conditional estimate and unique causal identification of legal institutions are different claims.

For the identification part, declare an admissible class \(\Theta\) of structural models, including preferences, technologies, shock laws, measurement/sampling rules and any equilibrium selector. If \(O\) denotes the complete observable history and \(T(\theta)\) the target defined above, the identified set is \(\mathcal I_T=\{T(\theta):\theta\in\Theta,\ \mathcal L_\theta(O)=\mathcal L_{\rm obs}(O)\}\); point identification means this set is a singleton. A \(do(a)\) comparison replaces stated policy/technology equations while fixing the declared exogenous law and re-solving the remaining equations. These definitions do not assert that an identification theorem holds without further restrictions.
\subsection*{Variants and assumption changes}
\begin{enumerate}
\item \textbf{Fixed contracting model (source subcase).} Estimate its declared enforcement/appropriability wedge and compute the model-specific welfare and trade counterfactuals; this conditional task already has published working-paper estimates.
\item \textbf{Institutional mechanism identification (editorial).} Separate enforcement, nonverifiability, search and finance using independently justified variation, allowing endogenous ownership and supplier investment.
\end{enumerate}
\subsection*{References and source audit}
\begin{enumerate}
\item NBER indexed primary record verifies Chor and Ma, 2025 WP 34044 and DOI; no invented priority claim. Locator: Indexed abstract; exact full-text introduction locator unverified. Support: Conditional contracting welfare/trade study. Original author PDF and NBER direct PDF/page fetches failed; primary indexed abstract/metadata inspected. URL: \url{https://www.nber.org/papers/w34044}.
\end{enumerate}
\begin{enumerate}\raggedright
\item\hypertarget{definitions:source:30007631:1}{} Davin Chor; Lin Ma. 2025. \emph{The Aggregate Welfare and Trade Implications of Contracting Frictions in Global Sourcing}. Indexed abstract; exact full-text introduction locator unverified.
\url{https://www.nber.org/papers/w34044}.
DOI: \url{https://doi.org/10.3386/w34044}.
\end{enumerate}

\subsection*{Common conventions: Source mathematical conventions}

These conventions apply to each entry unless explicitly replaced.

\textbf{Models and identification.} A primitive model \(m\) specifies all probability laws, feasible choices, information, equilibrium and measurement rules used by its target. The admissible class \(\mathcal M\) is fixed independently of outcomes. If \(P_O(m)\) is its observable probability law and \(\tau(m)\) the target, its identified set at observable law \(P\) is
\[
 \mathcal I_\tau(P)=\{\tau(m):m\in\mathcal M,\ P_O(m)=P\}.
\]
Point identification means that this set is a singleton. An empty set signals incompatibility of data and assumptions. Coordinate bounds need not describe the joint set. Bounds are sharp only if every retained target is attainable or arbitrarily approachable by compatible models. Finite bounds require explicit restrictions on unobserved quantities; unrestricted confounding, hidden wealth, extrapolation or arbitrary equilibrium selection can yield uninformative sets.

\textbf{Causal interventions.} A potential outcome \(Y_i(p)\) is the outcome under a complete policy regime \(p\), including timing, financing, information and market spillovers. The target population and units are fixed before treatment. With interference, use the complete assignment vector or a justified exposure mapping; individual no-interference notation alone cannot identify an economy-wide intervention. A difference of mean potential outcomes does not require the cross-world joint distribution. Conditioning on post-treatment migration, survival, enrollment or employment changes the estimand and needs separate justification.

\textbf{Statistical statements.} An estimator, confidence set, forecast or test specifies a sampling law, model class and loss/coverage criterion. Uniform \((1-\alpha)\)-coverage means \(\inf_{m\in\mathcal M}\Pr_m\{\tau(m)\in C_n\}\geq1-\alpha\), or an explicitly asymptotic version. The whole parameter space trivially has coverage, so informativeness requires a stated width, risk or power objective. A forecasting improvement is relative to a named benchmark, information set and evaluation population.

\textbf{Equilibrium and optimization.} An equilibrium comprises feasible plans, optimality, beliefs where needed, budget constraints and all market/resource clearing. The primitives or a stated benchmark define each correspondence \(\mathcal E(p;m)\). Nonemptiness is assumed or proved, never inferred just from compact policy sets. A selected-equilibrium target uses a declared selection rule; a robust target quantifies over all permitted equilibria. Use suprema or infima unless attainment is established. An \(\varepsilon\)-optimal policy is feasible and within \(\varepsilon\) of the finite optimum value. Report an empty feasible set separately.

\textbf{Welfare and accounting.} Normative utility, interpersonal normalization, social weights, numeraire, discounting and the use of public receipts are primitives. Behavioral utility need not coincide with the welfare criterion. Household welfare includes ownership income and rents once; adding profits again to compensating variation generally double-counts them. A monetary equivalent is defined by an explicit transfer and price convention, with monotonicity and range conditions. Average effects, mechanism contributions, transfers and welfare losses are different targets. Infinite sums require convergence and dynamic feasible plans require terminal or no-Ponzi conditions.

\textbf{Source versions.} A working-paper date, revision date and journal publication year are distinguished. A DOI for a journal version is not automatically the DOI of an earlier manuscript. Locators refer to the stated version and printed pages where specified. A metadata-only check does not verify a claim about a full-text conclusion. Every entry's source subsection narrows the provenance claim accordingly.
% END ECONOMICS STATEMENT
\end{document}
